% ---------------------------------------------------
% Trabajo Final de Grado
% Author: Fabián González Lence <alu0101549491@ull.edu.es>
% Chapter: Conclusiones y Líneas de Trabajo Futuras
% ----------------------------------------------------

\chapter{Conclusiones y Líneas de Trabajo Futuras} \label{chap:Conclusiones}

\section{Conclusiones} \label{sec:conclusiones}

Fruto del trabajo realizado, el repositorio del \ac{TFG} contiene cinco aplicaciones web desplegadas y funcionales, cuyo código de producción se generó en su totalidad mediante modelos de \ac{IA} con roles diferenciados, capaces de producir en ocasiones, en una única iteración de \textit{prompt}, módulos completos de código estructurado, documentado y coherente con los diagramas de diseño. Poner esto en perspectiva es el punto de partida más honesto: en octubre de 2025, cuando se iniciaba el experimento con \ac{HANGMAN} y se evaluaba si Claude era capaz de generar un diagrama UML en Mermaid, la posibilidad de llegar hasta aquí no era evidente. Los resultados obtenidos, analizados en el Capítulo~\ref{chap:Results}, superan las expectativas iniciales en términos de escala y coherencia, aunque esta valoración positiva convive con episodios de rendimiento mediocre, particularmente en tareas iterativas ---ajuste de interfaz de usuario, métricas de calidad, entre otras--- donde los modelos tendían a requerir varias rondas de instrucciones para converger al resultado esperado.\\

Uno de los hallazgos metodológicos más relevantes del trabajo es que la distribución de roles entre modelos ha demostrado ser la decisión de diseño con mayor impacto en la calidad del resultado final. Tal como se documenta en el Capítulo~\ref{chap:Metodology}, esta distribución introduce una capa de abstracción que reproduce la dinámica de un equipo de desarrollo humano; sin embargo, su beneficio pleno no se materializó hasta la transición ---descrita en el Capítulo~\ref{chap:AIModels}--- de \textit{chatbots} conversacionales externos al \ac{IDE} a agentes personalizados de GitHub Copilot integrados directamente en el repositorio. El acceso directo al código base, sin la fricción de copiar y pegar entre herramientas externas, propició un grado de especialización y desacoplamiento que redujo notablemente la incoherencia intermodular observada en los proyectos anteriores al cambio. Este cambio no fue planificado, sino que surgió de la evidencia acumulada en los tres primeros proyectos: el enfoque basado en \textit{chatbots} conversacionales es eficaz en proyectos de pequeño alcance como \ac{HANGMAN} y \ac{PLAYER}, pero no escala con suficiente coherencia cuando el tamaño del sistema crece hasta la envergadura de \ac{CARTO} o \ac{TENNIS}, siendo \ac{BALATRO} el punto de inflexión donde sus limitaciones se hicieron más evidentes.\\

En lo relativo al rol humano en el proceso, la experiencia confirma que el valor diferencial del ingeniero no residió en la escritura de código, sino en la capacidad de coordinación entre agentes, la validación de los resultados intermedios y la traslación de las preferencias del cliente al lenguaje de \textit{prompt}. Esta función, análoga a la de un director técnico que define qué construir y a quién encargárselo, resultó imprescindible durante todo el proceso: sin esa dirección, los agentes carecen de la orientación necesaria para mantener la coherencia global del producto. La capacidad de ejercer este rol requiere, a su vez, un conocimiento técnico profundo del dominio del desarrollo de software, lo que refuerza la idea de que la \ac{IA} no elimina la necesidad del ingeniero, sino que transforma la naturaleza de su trabajo.\\

Respecto a los modelos empleados, y tal como se analiza en los Capítulos~\ref{chap:AIModels} y~\ref{chap:Results}, Claude ---específicamente en sus versiones Sonnet 4.5 y 4.6--- destacó de forma consistente por combinar una gran ventana de contexto con una alta capacidad para generar código modular y coherente con las directrices técnicas establecidas en los \textit{prompts}. Qwen mostró un rendimiento notable en los proyectos de menor escala, especialmente en la generación de conjuntos de pruebas para \ac{HANGMAN} y \ac{PLAYER}, pero su rendimiento se degradó de forma significativa en \ac{BALATRO}, donde la complejidad del modelo de dominio superó su capacidad para producir pruebas funcionales de manera estable. Esta degradación fue uno de los factores que precipitaron el cambio de enfoque metodológico.\\

En cuanto a la calidad del código generado, los cinco proyectos presentan en general una estructura modular alineada con los diagramas de clase producidos en la fase de diseño, con documentación acorde a los estándares de ingeniería de software. El aspecto que con mayor frecuencia presenta deficiencias es la coherencia intermodular: cuando el contexto del proyecto supera la ventana efectiva del modelo, se producen dependencias inconsistentes, métodos duplicados o referencias a módulos inexistentes. Este fenómeno, especialmente acusado en \ac{BALATRO} durante la fase de codificación con Mistral, constituye uno de los argumentos más sólidos a favor del enfoque basado en agentes especializados con acceso directo al repositorio.\\

La reflexión de mayor alcance que emerge de este trabajo concierne al lugar que ocupan los agentes de \ac{IA} en el proceso de desarrollo de software. A partir de la experiencia acumulada a lo largo de los cinco proyectos, se concluye que, en el estado actual de la tecnología, estos modelos no reemplazan al ingeniero: actúan como una extensión del mismo, de manera análoga a como la aparición de los lenguajes de alto nivel no eliminó la figura del programador sino que transformó su forma de trabajar. La clave diferencial no reside en ser capaz de describir a una \ac{IA} lo que se desea construir ---eso está al alcance de cualquier persona---, sino en poseer el conocimiento técnico suficiente para evaluar la calidad de lo que se recibe, identificar sus deficiencias y dirigir el proceso de forma que el resultado final sea correcto, mantenible y escalable. Los resultados de este \ac{TFG} apuntan, en definitiva, a que el paradigma de desarrollo asistido por \ac{IA} no elimina la necesidad de la formación técnica en ingeniería del software, sino que eleva el umbral de lo que un ingeniero bien formado puede producir en un tiempo dado.

\section{Líneas de Trabajo Futuro} \label{sec:trabajo-futuro}

La escala del experimento, limitada a cinco proyectos web, condiciona la generalidad de las conclusiones obtenidas. Una primera línea de investigación futura sería aumentar el número de participantes y replicar el experimento con distintos perfiles de programador: pedir a varios ingenieros que construyan la misma aplicación empleando los mismos modelos y comparar los resultados permitiría aislar y estudiar el efecto de la habilidad de \textit{prompting} del experimentador del rendimiento intrínseco del modelo.\\

Complementariamente, retomar el enfoque comparativo descartado en las fases iniciales del trabajo ---desarrollar la misma aplicación de forma paralela con distintos modelos y comparar los resultados en términos de calidad, cobertura de requisitos y esfuerzo de supervisión--- proporcionaría datos cuantitativos sobre las diferencias entre modelos que el enfoque colaborativo adoptado no permite extraer de forma aislada, y respondería preguntas que este \ac{TFG} planteó pero no pudo responder empíricamente.\\

Extender el experimento a dominios distintos del desarrollo web constituye otra dirección relevante: aplicaciones de escritorio, sistemas embebidos o \textit{backends} de procesamiento intensivo de datos introducirían restricciones técnicas diferentes que pondrían a prueba las capacidades de los modelos en contextos menos representados en sus datos de entrenamiento y con otros lenguajes de programación.\\

Finalmente, la línea de mayor potencial transformador a largo plazo es la automatización completa del flujo de coordinación entre agentes: que una \ac{IA} gestione a otras \ac{IAs} sin intervención humana en la dirección del proceso. En el estado actual de la tecnología, el rol del jefe de desarrollo humano continúa siendo la pieza imprescindible del sistema; su sustitución por un agente orquestador autónomo representaría un avance cualitativo que, aunque técnicamente plausible, aún no ha demostrado ser viable en proyectos de complejidad real. Su investigación en el contexto del desarrollo de software asistido por \ac{IA} constituiría una contribución de alto impacto en el ámbito de la ingeniería del software.\\